\section*{A2.7. Línea base y respaldo del dimensionamiento}
Este anexo separa datos operacionales del caso, cálculos derivados y escenarios de diseño. Los volúmenes operacionales no equivalen a usuarios concurrentes ni transacciones informáticas. Las proyecciones conservan su condición de referencia o ampliación en evaluación; no se presentan como crecimiento contratado.

\subsection*{Volumetría operacional}
La fuente de las veinte dimensiones es C07, numeral 14.1. Se conserva la aproximación de la fuente. Su columna de proyección se titula «a tres años»; el escenario de expansión física se vincula además con 2031–2033 en RT-02.12, sin forzar equivalencia de fechas.

La Tabla~\ref{tab:a2-7-01} detalla las magnitudes y el respaldo del dimensionamiento.
\setlength{\AnchoTablaAnexo}{\dimexpr\linewidth-10\tabcolsep-6\arrayrulewidth\relax}
\begin{longtable}{|L{0.075697\AnchoTablaAnexo}|L{0.254980\AnchoTablaAnexo}|R{0.294821\AnchoTablaAnexo}|R{0.266932\AnchoTablaAnexo}|L{0.107570\AnchoTablaAnexo}|}
\caption{Volumetría operacional}\label{tab:a2-7-01}\\
\hline
\EncabezadoTabla{ID} & \EncabezadoTabla{Dimensión} & \EncabezadoTabla{Valor actual} & \EncabezadoTabla{Proyección del caso} & \EncabezadoTabla{Tipo / fuente} \\ \hline
\endfirsthead
\multicolumn{5}{l}{\small\textit{Línea base — continuación}}\\
\hline
\EncabezadoTabla{ID} & \EncabezadoTabla{Dimensión} & \EncabezadoTabla{Valor actual} & \EncabezadoTabla{Proyección del caso} & \EncabezadoTabla{Tipo / fuente} \\ \hline
\endhead
\multicolumn{5}{r}{\small\textit{}}\\
\endfoot
\endlastfoot
\rowcolor{RFclaro}
DIM-01 & Amarras en el agua & 320 & hasta 380, si se ejecuta la ampliación en evaluación & Dato / C07 sec. 14.1 \\ \hline
DIM-02 & Contratos permanentes de amarra vigentes & 296 & ≈ 350 & Dato / C07 sec. 14.1 \\ \hline
\rowcolor{RFclaro}
DIM-03 & Puestos de invernada en seco & 180, con 152 ocupados & 200 & Dato / C07 sec. 14.1 \\ \hline
DIM-04 & Boyas de fondeo en Bahía Ilque & 40 & hasta 70 & Dato / C07 sec. 14.1 \\ \hline
\rowcolor{RFclaro}
DIM-05 & Recaladas de embarcaciones visitantes al año & ≈ 1.900 & ≈ 2.400 & Dato / C07 sec. 14.1 \\ \hline
DIM-06 & Zarpes tramitados al año & ≈ 4.800 & ≈ 5.600 & Dato / C07 sec. 14.1 \\ \hline
\rowcolor{RFclaro}
DIM-07 & Zarpes en el peak de un fin de semana largo & hasta 90, concentrados en dos mañanas & hasta 110 & Dato / C07 sec. 14.1 \\ \hline
DIM-08 & Movimientos de travelift al año & 1.150 & 1.350 & Dato / C07 sec. 14.1 \\ \hline
\rowcolor{RFclaro}
DIM-09 & Litros de combustible expendidos al año & 1.900.000 & 2.200.000 & Dato / C07 sec. 14.1 \\ \hline
DIM-10 & Socios del club & 1.240 & ≈ 1.450 & Dato / C07 sec. 14.1 \\ \hline
\rowcolor{RFclaro}
DIM-11 & Alumnos de la escuela de vela al año & 640, de ellos 470 menores de edad & 800, de ellos ≈ 590 menores & Dato / C07 sec. 14.1 \\ \hline
DIM-12 & Salidas de clase al agua al año & ≈ 1.100, con hasta 18 embarcaciones menores cada una & ≈ 1.400 & Dato / C07 sec. 14.1 \\ \hline
\rowcolor{RFclaro}
DIM-13 & Regatas y eventos náuticos al año & 14; el mayor con 120 embarcaciones y 800 personas & 18 & Dato / C07 sec. 14.1 \\ \hline
DIM-14 & Empresas contratistas registradas & 62, con ≈ 140 personas & 75, con ≈ 170 personas & Dato / C07 sec. 14.1 \\ \hline
\rowcolor{RFclaro}
DIM-15 & Puntos de medición de agua y energía a instalar & 0 hoy; 320 amarras x 2 servicios & hasta 760 puntos con la ampliación & Dato / C07 sec. 14.1 \\ \hline
DIM-16 & Torretas de servicio en pantalán & 168 & 200 & Dato / C07 sec. 14.1 \\ \hline
\rowcolor{RFclaro}
DIM-17 & Cámaras de videovigilancia & 38 & 60 & Dato / C07 sec. 14.1 \\ \hline
DIM-18 & Documentos de cobro emitidos al mes & ≈ 1.950 & ≈ 2.300 & Dato / C07 sec. 14.1 \\ \hline
\rowcolor{RFclaro}
DIM-19 & Personal propio & 39 en invierno, hasta 78 en peak & 44 y 88 & Dato / C07 sec. 14.1 \\ \hline
DIM-20 & Personas distintas que ingresan al recinto en un mes de temporada & ≈ 3.400 entre socios, tripulantes, apoderados, contratistas y visitas & ≈ 4.000 & Dato / C07 sec. 14.1 \\ \hline
\end{longtable}
\FuenteTabla{Registro de Synaptix; fuentes y vínculos identificados en las filas.}

Las fuentes y los supuestos mantienen visible la diferencia entre datos de base, cálculos y escenarios.

\subsection*{Cálculos derivados}
Los resultados siguientes se obtienen de los datos citados; los redondeos porcentuales se expresan con dos decimales y los volúmenes aproximados conservan su carácter estimado.

La Tabla~\ref{tab:a2-7-02} detalla las magnitudes y el respaldo del dimensionamiento.
\setlength{\AnchoTablaAnexo}{\dimexpr\linewidth-10\tabcolsep-6\arrayrulewidth\relax}
\begin{longtable}{|L{0.176707\AnchoTablaAnexo}|L{0.204819\AnchoTablaAnexo}|R{0.148594\AnchoTablaAnexo}|L{0.116466\AnchoTablaAnexo}|L{0.353414\AnchoTablaAnexo}|}
\caption{Cálculos derivados}\label{tab:a2-7-02}\\
\hline
\EncabezadoTabla{ID / indicador} & \EncabezadoTabla{Fórmula} & \EncabezadoTabla{Resultado} & \EncabezadoTabla{Fuente} & \EncabezadoTabla{Interpretación y límite} \\ \hline
\endfirsthead
\multicolumn{5}{l}{\small\textit{Cálculos — continuación}}\\
\hline
\EncabezadoTabla{ID / indicador} & \EncabezadoTabla{Fórmula} & \EncabezadoTabla{Resultado} & \EncabezadoTabla{Fuente} & \EncabezadoTabla{Interpretación y límite} \\ \hline
\endhead
\multicolumn{5}{r}{\small\textit{}}\\
\endfoot
\endlastfoot
\rowcolor{RFclaro}
CAL-01 — Ocupación de invernada & 152 / 180 \(\times\) 100 & 84,44\% & C07 sec. 14.1 & Ocupación del inventario de invernada; no ocupación de amarras. \\ \hline
CAL-02 — Dotación estacional & 78 / 39 & 2 veces & C07 sec. 14.1 & Duplicación de personas; no duplicación automática de usuarios concurrentes. \\ \hline
\rowcolor{RFclaro}
CAL-03 — Movimientos en campañas & 1.150 \(\times\) 0,62 & 713 movimientos & C07 sec. 7.3 & Total conjunto de dos campañas; no reparto idéntico ni ritmo diario constante. \\ \hline
CAL-04 — Tramitación anual de zarpe & 4.800 \(\times\) 18 / 60 & ≈1.440 horas/año & C07 sec. 7.2 & Carga acumulada aproximada; no ahorro prometido ni duración de proyecto. \\ \hline
\rowcolor{RFclaro}
CAL-05 — Trámites con corrección & 4.800 \(\times\) 0,34 & ≈1.632 trámites/año & C07 sec. 7.2 & Estimación con tasa agregada; no sumar tiempo sin conocer si ya integra los 18 minutos. \\ \hline
CAL-06 — Contactados en marejada & 190 / 296 \(\times\) 100 & 64,19\% & C07 sec. 7.2 & Resultado del episodio; no tasa general de entrega de una aplicación. \\ \hline
\rowcolor{RFclaro}
CAL-07 — No contactados en episodio & 296 - 190 & 106 armadores & C07 sec. 7.2 & Complemento del recuento del aviso, no personas siempre incontactables. \\ \hline
CAL-08 — Ampliación de amarras & (380 - 320) / 320 \(\times\) 100 & 18,75\% & C07 RT-02.12 & Escenario físico en evaluación; no tasa de aumento de transacciones. \\ \hline
\rowcolor{RFclaro}
CAL-09 — Ampliación de boyas & (70 - 40) / 40 \(\times\) 100 & 75\% & C07 RT-02.12 & Capacidad futura condicionada a ampliación. \\ \hline
CAL-10 — Puntos de medición base & 320 \(\times\) 2 & 640 puntos & C07 sec. 14.1 & Escenario de dos servicios por amarra; equipos físicos dependen del diseño. \\ \hline
\rowcolor{RFclaro}
CAL-11 — Puntos con ampliación & 380 \(\times\) 2 & 760 puntos & C07 sec. 14.1 & No equivale a 760 torretas ni determina frecuencia de muestreo. \\ \hline
CAL-12 — Participación de menores & 470 / 640 \(\times\) 100 & 73,44\% & C07 sec. 14.1 & Fracción de matrícula anual; no alumnos simultáneamente en agua. \\ \hline
\rowcolor{RFclaro}
CAL-13 — Preparación de facturación & 10 días-persona/mes \(\times\) 12 & 120 días-persona/año & C07 sec. 7.4 & Esfuerzo equivalente; no plazo calendario ni jornada en horas no declaradas. \\ \hline
\end{longtable}
\FuenteTabla{Registro de Synaptix; fuentes y vínculos identificados en las filas.}

Las fuentes y los supuestos mantienen visible la diferencia entre datos de base, cálculos y escenarios.

\subsection*{Escenarios que no deben confundirse con promedios}
El escenario de C07 sec. 14.2 reúne 90 zarpes, 20 recaladas, tres clases en el agua y campaña de varada entre 07:00 y 11:00. Debe traducirse a operaciones, interacciones e integraciones antes de estimar carga. El dato de sec. 14.1 sobre dos mañanas se conserva con su propio contexto. Los 40 minutos del incidente de escuela son un episodio, no un promedio de atención ni un tiempo de exposición al peligro.
\subsection*{Control de estimaciones de sistema}
Las dieciocho dimensiones de C07 sec. 14.2 requieren valor y justificación del proponente. Las estimaciones de sistema se obtienen a partir de la volumetría del caso y de supuestos explícitos sobre frecuencia de uso, simultaneidad, tamaño de los registros y condiciones de comunicación. Son estimaciones de formulación, no mediciones de una solución implementada. Cada resultado identifica su unidad, escenario, método de cálculo y referencia al desarrollo que lo sustenta. Los ensayos posteriores contrastarán las hipótesis utilizadas.

La Tabla~\ref{tab:a2-7-03} detalla las magnitudes y el respaldo del dimensionamiento.
\setlength{\AnchoTablaAnexo}{\dimexpr\linewidth-8\tabcolsep-5\arrayrulewidth\relax}
\begin{longtable}{|L{0.075697\AnchoTablaAnexo}|L{0.314741\AnchoTablaAnexo}|L{0.434263\AnchoTablaAnexo}|L{0.175299\AnchoTablaAnexo}|}
\caption{Control de estimaciones de sistema}\label{tab:a2-7-03}\\
\hline
\EncabezadoTabla{ID} & \EncabezadoTabla{Dimensión a estimar} & \EncabezadoTabla{Método de obtención} & \EncabezadoTabla{Valor / referencia pendiente} \\ \hline
\endfirsthead
\multicolumn{4}{l}{\small\textit{Estimaciones de sistema — continuación}}\\
\hline
\EncabezadoTabla{ID} & \EncabezadoTabla{Dimensión a estimar} & \EncabezadoTabla{Método de obtención} & \EncabezadoTabla{Valor / referencia pendiente} \\ \hline
\endhead
\multicolumn{4}{r}{\small\textit{}}\\
\endfoot
\endlastfoot
\rowcolor{RFclaro}
EST-01 & Transacciones por segundo en régimen normal & Sumar operaciones por flujo y tasa de llegada por intervalo; separar lectura, escritura y procesos automáticos. & FALTA ESTIMACIÓN y referencia a SD4/SD5 o modelo de servicio, con supuestos y prueba. \\ \hline
EST-02 & Transacciones por segundo en el peak de una mañana de viernes largo, con zarpes, visitantes y clases simultáneas & Modelar simultaneidad del escenario de sec. 14.2 con duración de cada flujo, ráfagas y reintentos; comprobar mediante carga. & FALTA ESTIMACIÓN y referencia a SD4/SD5 o modelo de servicio, con supuestos y prueba. \\ \hline
\rowcolor{RFclaro}
EST-03 & Eventos por segundo generados por la medición de agua y energía de las 320 amarras, según la frecuencia de muestreo que se proponga & Eventos/s = suma de puntos por servicio dividida por intervalo de muestreo en segundos; añadir cambios de ocupación y reintentos. & FALTA ESTIMACIÓN y referencia a SD4/SD5 o modelo de servicio, con supuestos y prueba. \\ \hline
EST-04 & Frecuencia de muestreo justificada para consumo eléctrico, consumo de agua y estado de ocupación de amarra & Definir y justificar el intervalo de adquisición, de transmisión y de actualización de la información, indicando unidades. Para agua y energía, comprobar la actualización al menos horaria del consumo acumulado por amarra. Para ocupación, comprobar el plazo específico aplicable desde el evento. Utilizar las frecuencias resultantes para calcular eventos, tráfico y almacenamiento. & FALTA ESTIMACIÓN y referencia a SD4/SD5 o modelo de servicio, con supuestos y prueba. \\ \hline
\rowcolor{RFclaro}
EST-05 & Personas usuarias internas concurrentes en un día de temporada & Estimar por perfil, turno y fracción simultánea de uso; no igualar concurrencia a dotación máxima. & FALTA ESTIMACIÓN y referencia a SD4/SD5 o modelo de servicio, con supuestos y prueba. \\ \hline
EST-06 & Personas usuarias externas concurrentes en el portal, distinguiendo armadores, apoderados, visitantes y contratistas & Separar población, sesiones por período y duración por perfil externo; aplicar escenarios y validar sensibilidad. & FALTA ESTIMACIÓN y referencia a SD4/SD5 o modelo de servicio, con supuestos y prueba. \\ \hline
\rowcolor{RFclaro}
EST-07 & Volumen anual de almacenamiento transaccional & Registros/año por tipo \(\times\) bytes por registro, más índices, auditoría y crecimiento; declarar factores y medir una muestra. & FALTA ESTIMACIÓN y referencia a SD4/SD5 o modelo de servicio, con supuestos y prueba. \\ \hline
EST-08 & Volumen anual de almacenamiento de las series de consumo por amarra & Puntos \(\times\) muestras/día \(\times\) días/año \(\times\) bytes/muestra; incluir metadatos, retención y compresión acreditada. & FALTA ESTIMACIÓN y referencia a SD4/SD5 o modelo de servicio, con supuestos y prueba. \\ \hline
\rowcolor{RFclaro}
EST-09 & Volumen de la evidencia documental digitalizada: contratos, matrículas, seguros, títulos, fichas de escuela y actas ambientales & Documentos por tipo \(\times\) tamaño medio de muestra \(\times\) versiones y años conservados; separar originales y miniaturas. & FALTA ESTIMACIÓN y referencia a SD4/SD5 o modelo de servicio, con supuestos y prueba. \\ \hline
EST-10 & Volumen total de datos históricos a migrar desde el sistema de 2014 & Perfilar y medir fuentes reales por tabla/documento; distinguir alcance migrado de archivo histórico. & FALTA ESTIMACIÓN y referencia a SD4/SD5 o modelo de servicio, con supuestos y prueba. \\ \hline
\rowcolor{RFclaro}
EST-11 & Número de integraciones y volumen de mensajes por integración & Inventariar interfaces efectivas y calcular mensajes por evento, respuestas, confirmaciones y reintentos. & FALTA ESTIMACIÓN y referencia a SD4/SD5 o modelo de servicio, con supuestos y prueba. \\ \hline
EST-12 & Cobertura y ancho de banda necesarios en los cinco pantalanes, con estructura en movimiento & Dimensionar por flujos simultáneos, telemetría y archivos; medir cobertura y caudal en marea y flexión. & FALTA ESTIMACIÓN y referencia a SD4/SD5 o modelo de servicio, con supuestos y prueba. \\ \hline
\rowcolor{RFclaro}
EST-13 & Tamaño máximo, frecuencia y latencia tolerable de los mensajes intercambiados con una embarcación en navegación por enlace de baja capacidad & Definir payload mínimo, frecuencia, latencia de horas y tolerancia a desorden según canal verificado. & FALTA ESTIMACIÓN y referencia a SD4/SD5 o modelo de servicio, con supuestos y prueba. \\ \hline
EST-14 & Volumen de datos generado por el recinto durante 24 horas sin enlace hacia el exterior & Sumar registros y adjuntos de todos los procesos offline durante 24 horas más margen explícito y validado. & FALTA ESTIMACIÓN y referencia a SD4/SD5 o modelo de servicio, con supuestos y prueba. \\ \hline
\rowcolor{RFclaro}
EST-15 & Tiempo de sincronización tras 24 horas de operación desconectada & Tiempo = datos a transferir / caudal útil más conciliación; probar cola, reintentos y límite de 30 minutos. & FALTA ESTIMACIÓN y referencia a SD4/SD5 o modelo de servicio, con supuestos y prueba. \\ \hline
EST-16 & Contactos mensuales a la mesa de ayuda, en español e inglés, distinguiendo personal interno de clientes externos & Proyectar contactos por perfil, idioma, estacionalidad y tipo de incidente; declarar duración de atención. & FALTA ESTIMACIÓN y referencia a SD4/SD5 o modelo de servicio, con supuestos y prueba. \\ \hline
\rowcolor{RFclaro}
EST-17 & Dotación de la mesa de ayuda y del equipo de operación, considerando que el CLIENTE no aporta ninguna persona de tecnologías de información & Calcular turnos y ocupación con modelo de colas/Erlang C y SLA, descansos, ausencias y cobertura de críticos. & FALTA ESTIMACIÓN y referencia a SD4/SD5 o modelo de servicio, con supuestos y prueba. \\ \hline
EST-18 & Esfuerzo mensual del rol de administración funcional de la solución y a quién se asigna & Sumar tareas \(\times\) frecuencia \(\times\) duración, gestión de cambios y suplencias; asignar rol y costear 36 meses. & FALTA ESTIMACIÓN y referencia a SD4/SD5 o modelo de servicio, con supuestos y prueba. \\ \hline
\end{longtable}
\FuenteTabla{Registro de Synaptix; fuentes y vínculos identificados en las filas.}

Las fuentes y los supuestos mantienen visible la diferencia entre datos de base, cálculos y escenarios.

La frecuencia horaria de actualización de consumos no impone una lectura cada hora: la frecuencia de adquisición se debe justificar por separado. Tampoco corresponde traducir el 38\% de costos de servicios no recuperados a volumen físico desperdiciado. El control VAC-15 registra la corrección documental del recuento a las tres no conformidades ambientales indicadas en el documento del Caso 07. Esa corrección no representa el cierre de los hallazgos ambientales.
